\documentclass[11pt]{article}
\usepackage{color}
%\input{rgb}

%----------Packages----------
\usepackage{amsmath}
\usepackage{amssymb}
\usepackage{amsthm}
\usepackage{amsrefs}
%\usepackage{dsfont}
\usepackage{mathrsfs}
%\usepackage{stmaryrd}
%\usepackage[all]{xy}
\usepackage[mathcal]{eucal} % changes meaning of \mathcal
\usepackage{enumerate}
\usepackage[shortlabels]{enumitem}
\usepackage{verbatim}  %%includes comment environment
\usepackage{fullpage}  %%smaller margins

%----------Commands----------
%%penalizes orphans
\clubpenalty=9999
\widowpenalty=9999

%% blackboard bold math capitals
\newcommand{\bbF}{\mathbb{F}}
\newcommand{\bbN}{\mathbb{N}}
\newcommand{\bbQ}{\mathbb{Q}}
\newcommand{\bbR}{\mathbb{R}}
\newcommand{\bbZ}{\mathbb{Z}}

\renewcommand{\phi}{\varphi}
\renewcommand{\emptyset}{\O}

\newcommand{\abs}[1]{\lvert #1 \rvert}
\newcommand{\norm}[1]{\lVert #1 \rVert}

\newcommand{\sarr}{\rightarrow}
\newcommand{\arr}{\longrightarrow}

\DeclareMathOperator{\ext}{ext}
\DeclareMathOperator{\ho}{hole}

%----------Theorems----------

\newtheorem{theorem}{Theorem}[section]
\newtheorem{proposition}[theorem]{Proposition}
\newtheorem{lemma}[theorem]{Lemma}
\newtheorem{corollary}[theorem]{Corollary}
\newtheorem{examples}[theorem]{Examples}
\newtheorem{example}[theorem]{Example}

\newtheorem{axiom}{Axiom}

\theoremstyle{definition}
\newtheorem{definition}[theorem]{Definition}
\newtheorem*{definition*}{Definition}
\newtheorem{nondefinition}[theorem]{Non-Definition}
\newtheorem{exercise}[theorem]{Exercise}
\newtheorem{remark}[theorem]{Remark}
\newtheorem{warning}[theorem]{Warning}
\newtheorem{question}[theorem]{Question}


\newcommand{\head}[1]{
	\begin{center}
		{\large #1}
		\vspace{.2 in}
	\end{center}
	
	\bigskip
}
\newcommand{\hint}[2][Hint]{
	
	(#1: #2)
}

\numberwithin{equation}{subsection}

%----------Editting----------

%\newcommand{\hide}[1]{} % for student version
%\newcommand{\com}[1]{} % for student version
%\newcommand{\meta}[1]{} % for removing meta comments in the script

\newcommand{\hide}[1]{{\color{red} #1}} % for instructor version
\newcommand{\com}[1]{{\color{blue} #1}} % for instructor version
\newcommand{\meta}[1]{{\color{green} #1}} % for making notes about the script that are not intended to end up in the script

\begin{document}

\head{MATH 162, Winter 2025\\
SCRIPT 7: The Field Axioms} 



%%---  sheet number for theorem counter
\setcounter{section}{7}   

We will formalize the notions of addition and multiplication in structures called fields.  A field with a compatible order is called an ordered field. We will see that $\bbQ$ and $\bbR$ are both examples of ordered fields.

\begin{definition} 
	A \emph{binary operation} on a nonempty set $X$ is a function 
	\begin{align*}
		f \colon X \times X \arr X.
	\end{align*}
	We say that $f$ is \emph{associative} if
	\begin{align*}
		f(f(x, y), z) = f(x, f(y, z)) \quad \text{for all $x, y, z \in X$.}
	\end{align*}
	We say that $f$ is \emph{commutative} if
	\begin{align*}
		f(x, y) = f(y, x) \quad \text{for all $x, y \in X$.}
	\end{align*}
	An \emph{identity element} of a binary operation $f$ is an element $e \in X$ such that
	\begin{align*}
		f(x, e) = f(e, x) = x \quad \text{for all $x \in X$.}
	\end{align*}
\end{definition}

\begin{remark}
	Frequently, we denote a binary operation differently. If $*\colon X\times X\arr X$ is the binary operation, we often write $a*b$ in place of $*(a,b)$. We sometimes indicate this same operation by writing $(a,b)\mapsto a*b$.
\end{remark}

\begin{exercise}
	Rewrite Definition 7.1 using the notation of Remark 7.2.
\begin{proof}


\renewcommand\qedsymbol{QED}
\end{proof}
\end{exercise}

\begin{examples}  \hspace{1in}
	\begin{enumerate}
		\item  The function $+ \colon \bbZ \times \bbZ \arr \bbZ$ which sends a pair of integers $(m, n)$ to $+(m, n) = m + n$ is a binary operation on the integers, called addition.  Addition is associative, commutative and has identity element $0$.
		
		\item  The maximum of $m$ and $n$, denoted $\max(m, n)$, is an associative and commutative binary operation on $\bbZ$.  Is there an identity element for $\max$? 

        No.
		
		\item  Let $\wp(Y)$ be the power set of a set $Y$.  Recall that the power set consists of all subsets of $Y$.  Then the intersection of sets, $(A, B) \mapsto A \cap B$, defines an associative and commutative binary operation on $\wp(Y)$.  Is there an identity element for $\cap$?

Yes. $Y$ is an identity element for $\cap$.

        
	\end{enumerate}
\end{examples}

\begin{exercise}  
	Find a binary operation on a set that is not commutative.  
	Find a binary operation on a set that is not associative.
    
\begin{proof}


\renewcommand\qedsymbol{QED}
\end{proof}

\end{exercise}

\begin{exercise}
	Let $X$ be a nonempty finite set, and let $Y=\{f\colon X \arr X \mid \mbox{$f$ is bijective}\}$.  Consider the
	binary operation of composition of functions, denoted $\circ\colon Y \times Y \arr Y$ and defined by
	$(f \circ g)(x) = f(g(x))$, as seen in Definition~1.27.  Decide whether or not composition is commutative and/or associative and whether or not it has an identity.

\begin{proof}
For a counterexample, consider $X = \{1,2,3\}$. Define $f\colon X \rightarrow X$ as $f(1)=2,f(2)=3,f(3)=1$ and define $g\colon X \rightarrow X$ as $g(1)=3,g(2)=2,g(3)=1$. Then, $f(g(1))=1$ and $g(f(1))=2$. Hence, composition is not necessarily commutative (at least for $|X|\geq 3$).

To prove that composition is associative, it will suffice to show that $((f\circ g)\circ h)(x)=(f\circ(g\circ h))(x)$. Consider the following algebra.
        \begin{align*}
            ((f\circ g)\circ h)(x) &= (f\circ g)(h(x))\\
            &= f(g(h(x)))\\
            &= f((g\circ h)(x))\\
            &= (f\circ(g\circ h))(x)
        \end{align*}\par

Let $i \colon X \rightarrow X$ such that $i(x)=x$. Then, we want to show that for all $f\in Y$, $f\circ i=i\circ f=f$. Consider the following algebra.
        \begin{align*}
            (f\circ i)(x) &= f(i(x))\\
            &= f(x)\\
            &= i(f(x))\\
            &= (i\circ f)(x)
        \end{align*}

\renewcommand\qedsymbol{QED}
\end{proof}

\end{exercise}



\begin{theorem}
	Identity elements are unique.  That is, suppose that $f$ is a binary operation on a set $X$ that has two identity elements $e$ and $e'$.  Then $e = e'$.

\begin{proof}
$e = f(e, e') = e'$. 

\renewcommand\qedsymbol{QED}
\end{proof}

\end{theorem}



\begin{definition}
	A \emph{field} is a set $F$ with two binary operations  on $F$ called addition, denoted $+$, and multiplication, denoted $\cdot$\;, satisfying the following \emph{field axioms}:
	\begin{enumerate}[{FA}1]
		\item  (Commutativity of Addition)  For all $x,y\in F$, $x + y = y + x$.
		\item  (Associativity of Addition)  For all $x,y,z\in F$, $(x + y) + z = x + (y + z)$.
		\item  (Additive Identity)  There exists an element $0 \in F$ such that $x + 0 = 0 + x = x$ for all $x \in F$.
		\item  (Additive Inverses)  For any $x \in F$, there exists $y \in F$ such that $x + y = y + x = 0$,  called an additive inverse of $x$.
		\item  (Commutativity of Multiplication)  For all $x,y\in F$, $x \cdot y = y \cdot x$.
		\item  (Associativity of Multiplication)  For all $x, y, z \in F$, $(x \cdot y) \cdot z = x \cdot (y \cdot z)$.
		\item  (Multiplicative Identity)  There exists an element $1 \in F$ such that $x \cdot 1 = 1 \cdot x = x$ for all $x \in F$.
		\item  (Multiplicative Inverses)  For any $x \in F$ such that $x \neq 0$, there exists $y \in F$ such that $x \cdot y = y \cdot x = 1$,  called a multiplicative inverse of $x$.
		\item  (Distributivity of Multiplication over Addition)  For all $x, y, z \in F$,
		$x \cdot(y + z) = x \cdot y + x \cdot z$. 
		\item  (Distinct Additive and Multiplicative Identities)  $1 \neq 0$.
\end{enumerate}

\end{definition}

\begin{exercise}
	Consider the set $\bbF_{2} = \{0, 1\}$, and define binary operations $+$ and $\cdot$ on $\bbF_{2}$ by:
	\begin{center}
	$
	\begin{array}{ccccccc}
		0 + 0 = 0 & \phantom{MM} & 0 + 1 = 1 & \phantom{MM} & 1 + 0 = 1 & \phantom{MM} &1 + 1 = 0 \\
		0 \cdot 0 = 0 & \phantom{MM} & 0 \cdot 1 = 0 & \phantom{MM}  & 1\cdot 0 =0 & \phantom{MM} &1 \cdot 1 = 1 
	\end{array}
	$
	\end{center}
	
	Show that $\bbF_{2}$ is a field.  


\begin{proof}
        To prove that $\bbF_2$ obeys FA1, it will suffice to show that $0+0=0+0$, $0+1=1+0$, and $1+1=1+1$. The first and third of these are evidently true. For the second, we have $0+1=1=1+0$, so it is good, too.\par
        To prove that $\bbF_2$ obeys FA2, the following casework will suffice.
        \begin{align*}
            (0+0)+0 &= 0 = 0+(0+0)&
            (0+0)+1 &= 1 = 0+(0+1)\\
            (0+1)+0 &= 1 = 0+(1+0)&
            (1+0)+0 &= 1 = 1+(0+0)\\
            (0+1)+1 &= 0 = 0+(1+1)&
            (1+1)+0 &= 0 = 1+(1+0)\\
            (1+0)+1 &= 0 = 1+(0+1)&
            (1+1)+1 &= 1 = 1+(1+1)
        \end{align*}
        To prove that $\bbF_2$ obeys FA3, it will suffice to find an element $0\in\bbF_2$ such that $x+0=0+x=x$. Since $0+0=0$, $1+0=0$, and by commutativity, it is clear that 0 is an additive identity in $\bbF_2$.\par
        To prove that $\bbF_2$ obeys FA4, it will suffice to show that for all $x\in\bbF_2$, there exists a $y\in\bbF_2$ such that $x+y=y+x=0$. For 0, this object is 0 (since $0+0=0+0=0$), and for 1, this object is 1 (since $1+1=1+1=0$).\par
        To prove that $\bbF_2$ obeys FA5, it will suffice to show that $0\cdot 0=0\cdot 0$, $0\cdot 1=1\cdot 0$, and $1\cdot 1=1\cdot 1$. The first and third of these are evidently true. For the second, we have $0\cdot 1=0=1\cdot 0$, so it is good, too.\par
        To prove that $\bbF_2$ obeys FA6, the following casework will suffice.
        \begin{align*}
            (0\cdot 0)\cdot 0 &= 0 = 0\cdot (0\cdot 0)&
            (0\cdot 0)\cdot 1 &= 0 = 0\cdot (0\cdot 1)\\
            (0\cdot 1)\cdot 0 &= 0 = 0\cdot (1\cdot 0)&
            (1\cdot 0)\cdot 0 &= 0 = 1\cdot (0\cdot 0)\\
            (0\cdot 1)\cdot 1 &= 0 = 0\cdot (1\cdot 1)&
            (1\cdot 1)\cdot 0 &= 0 = 1\cdot (1\cdot 0)\\
            (1\cdot 0)\cdot 1 &= 0 = 1\cdot (0\cdot 1)&
            (1\cdot 1)\cdot 1 &= 1 = 1\cdot (1\cdot 1)
        \end{align*}
        To prove that $\bbF_2$ obeys FA7, it will suffice to find an element $1\in\bbF_2$ such that $x\cdot 1=1\cdot x=x$. Since $0\cdot 1=0$, $1\cdot 1=1$, and by commutativity, it is clear that 1 is a multiplicative identity in $\bbF_2$.\par
        To prove that $\bbF_2$ obeys FA8, it will suffice to show that for all $x\in\bbF_2$ such that $x\neq 0$, there exists a $y\in\bbF_2$ such that $x\cdot y=y\cdot x=1$. For 1, this object is 1 (since $1\cdot 1=1\cdot 1=1$).\par
        To prove that $\bbF_2$ obeys FA9, the following casework will suffice.
        \begin{align*}
            0\cdot(0+0) &= 0 = 0\cdot 0+0\cdot 0&
            0\cdot(0+1) &= 0 = 0\cdot 0+0\cdot 1\\
            0\cdot(1+0) &= 0 = 0\cdot 1+0\cdot 0&
            1\cdot(0+0) &= 0 = 1\cdot 0+1\cdot 0\\
            0\cdot(1+1) &= 0 = 0\cdot 1+0\cdot 1&
            1\cdot(1+0) &= 1 = 1\cdot 1+1\cdot 0\\
            1\cdot(0+1) &= 1 = 1\cdot 0+1\cdot 1&
            1\cdot(1+1) &= 0 = 1\cdot 1+1\cdot 1
        \end{align*}
        To prove that $\bbF_2$ obeys FA10, it will suffice to show that $0\neq 1$. Clearly this is true.

This completes the proof.

\renewcommand\qedsymbol{QED}
\end{proof}

\end{exercise}



\begin{theorem}
	Suppose that $F$ is a field.  Then additive inverses are unique.  This means:

 Let $x \in F$.  If $y, y' \in F$ satisfy $x + y = 0$ and $x + y' = 0$, then $y = y'$.
\begin{proof}
$y=y+0=y+(x+y')=(y+x)+y'=(x+y)+y'=0+y'=y'$.

\renewcommand\qedsymbol{QED}
\end{proof}

\end{theorem}


We usually write $-x$ for the additive inverse of $x$.

\begin{corollary}
	If $x\in F$, then $-(-x)=x$.

\begin{proof}
$(-x)+x=0$ and $(-x)+(-(-x))=0$, hence $x=-(-x)$. 

\renewcommand\qedsymbol{QED}
\end{proof}

\end{corollary}



\begin{theorem}
	Let $F$ be a field, and let $a,b,c\in F$. If $a + b = a + c$, then $b = c$.
    
\begin{proof}
$a+b=a+c$ implies $(-a)+(a+b)=(-a)+(a+c)$, hence $((-a)+a)+b=((-a)+a)+c$, hence $0+b=0=c$, which implies $b=c$.

\renewcommand\qedsymbol{QED}
\end{proof}

\end{theorem}



\begin{theorem}
	Let $F$ be a field. If $a\in F$, then $a \cdot 0 = 0$.
    
\begin{proof}
We claim that if $x+x=x$, then $x=0$. To see this, consider $(-x)+(x+x)=(-x)+x=(-x+x)+x=(-x+x)$. Hence, $x=0$.

$a\cdot0+a\cdot0=a\cdot(0+0)=a\cdot0$. Hence, by the claim, $a\cdot0=0$.

\renewcommand\qedsymbol{QED}
\end{proof}


\end{theorem}




\begin{theorem}
	Suppose that $F$ is a field.  Then multiplicative inverses are unique.  This means: 
	
	Let $x \in F$.  If $y, y' \in F$ satisfy $x \cdot y = 1$ and $x \cdot y' = 1$, then $y = y'$.
    
\begin{proof}
$x\cdot y=1=x\cdot y'$. Then, $y\cdot (x+y)=y\cdot (x+y')$. Then, $(y \cdot x)\cdot y=(y \cdot x)\cdot y'$. Then, $(x \cdot y)\cdot y=(x \cdot y)\cdot y'$. Then, $1\cdot y=1\cdot y'$. Hence, $y=y'$.

\renewcommand\qedsymbol{QED}
\end{proof}


\end{theorem}

We usually write $x^{-1}$ or $\frac{1}{x}$ for the multiplicative inverse of $x$.

\begin{corollary}\com{ Note that you need 7.13 to justify the existence of $(x^{-1})^{-1}.$}

	If $x\in F$ and $x\neq 0$, then $(x^{-1})^{-1}=x$.
    
\begin{proof}
First, we claim that $x^{-1}=0$. Suppose $x^{-1}\neq 0$. Then, $x \cdot x^{-1} = 0$, and this is a contradiction. 

Hence, $(x^{-1})^{-1}\cdot x^{-1}=1$. Hence, $x \cdot x^{-1}=1$. Hence, $x=(x^{-1})^{-1}$.

\renewcommand\qedsymbol{QED}
\end{proof}


\end{corollary}


\begin{theorem}
	Let $F$ be a field, and let $a, b, c \in F$. If $a \cdot b = a \cdot c$ and $a \neq 0$, then $b = c$.
    
\begin{proof}
$\frac{1}{a}\cdot a \cdot b= \frac{1}{a}\cdot a \cdot c$. Hence, $1 \cdot b=1 \cdot c$. Hence, $b=c$.



\renewcommand\qedsymbol{QED}
\end{proof}


\end{theorem}


\begin{theorem}
	Let $F$ be a field, and let $a, b\in F$.  If $a \cdot b = 0$, then $a = 0$ or $b = 0$.
    
\begin{proof}
Suppose $a \neq 0$. Then, $a\cdot b=0$, hence $\frac{1}{a}\cdot a\cdot b=\frac{1}{a}\cdot 0$. Hence, $1\cdot b=\frac{1}{a}\cdot 0$, hence $1 \cdot b=0$, hence $b=0$.

Thus, if one of the elements is non-zero, then the other must be zero.


\renewcommand\qedsymbol{QED}
\end{proof}


\end{theorem}

\begin{lemma}
	Let $F$ be a field. If $a\in F,$ then $-a=(-1)a.$
    
\begin{proof}
$0=0\cdot a=((-1)+1)\cdot a=(-1)\cdot a + 1 \cdot a=(-1)\cdot a + a$.

Hence $(-1)\cdot a + a = 0$, which implies that $-a=(-1)\cdot a$.

\renewcommand\qedsymbol{QED}
\end{proof}


\end{lemma}

\begin{lemma}
	Let $F$ be a field. If $a, b\in F$, then $a \cdot (-b )= -(a \cdot b) = (-a) \cdot b$.

        
\begin{proof}
$\underline{a\cdot(-b)}=a\cdot((-1)\cdot b)=a\cdot(b\cdot(-1))=(a\cdot b)\cdot(-1)=(-1)\cdot(a \cdot b)=\underline{-(a\cdot b)}=(-1)\cdot(a \cdot b)=((-1)\cdot (a))\cdot b=\underline{(-a)\cdot b}$. 

\renewcommand\qedsymbol{QED}
\end{proof}


\end{lemma}

\begin{lemma}
	Let $F$ be a field. If $a, b\in F$, then $a\cdot b=(-a)\cdot (-b)$.
        
\begin{proof}
$(-a)\cdot (-b)=-(a \cdot (-b)) = -(-(a \cdot b))=a \cdot b$.

\renewcommand\qedsymbol{QED}
\end{proof}


\end{lemma}

Next, we discuss the notion of an ordered field.

\begin{definition}
	An \emph{ordered field} is a field $F$ equipped with an ordering $<$ (satisfying Definition 3.1) such that also:
	\begin{itemize}
		\item  Addition respects the ordering: if $x < y$, then $x + z < y + z$ for all $z \in F$.
		\item  Multiplication respects the ordering: if $0 < x$ and $0 < y$, then $0 < x \cdot y$.
	\end{itemize}
\end{definition}

\begin{definition}
	Suppose $F$ is an ordered field and $x\in F$.  If $0 < x$, we say that $x$ is \emph{positive}. If $x < 0$, we say that $x$ is \emph{negative}.
\end{definition}

For the remaining theorems, let $F$ be an ordered field.

\begin{lemma}
	If $0 < x$, then $-x < 0$.  Similarly, if $x < 0$, then $0 < -x$.
        
\begin{proof}
If $0<x$, then $0+(-x)<x+(-x)$, hence $-x<0$.

If $x<0$, then $x+(-x)<0+(-x)$, hence $0<-x$. 

\renewcommand\qedsymbol{QED}
\end{proof}


\end{lemma}


\begin{lemma}
	Let $x,y,z \in F$. 
	\begin{enumerate}
		\item If $x > 0$ and $y < z$, then $x \cdot y < x \cdot z$.
\begin{proof}
$y+(-y)<z+(-y)$, hence $0<z+(-y)$. Thus, $x\cdot (z+(-y))>0$, then $x\cdot z + x \cdot(-y) > 0$. Hence, $x \cdot z+(-(x \cdot y))>0$, hence $x\cdot z > x \cdot y$. 

\renewcommand\qedsymbol{QED}
\end{proof}
		\item If $x < 0$ and $y < z$, then $x \cdot z < x \cdot y$.
\begin{proof}
Applying item 1, we get that $(-x)\cdot y < (-x) \cdot z$. Then, $0<(-(x \cdot z))+(x \cdot y)$. Hence, $x\cdot z<x \cdot y$.

\renewcommand\qedsymbol{QED}
\end{proof}
	\end{enumerate}
	
        



\end{lemma}

\begin{remark} An immediate consequence of this lemma is the fact that if $x$ and $y$ are both positive or both negative, their product is positive.
\end{remark}


\begin{lemma}
	If $x \in F$, then $0 \leq x^2$.  Moreover, if $x \neq 0$, then $0 < x^2$.
        
\begin{proof}
If $x = 0$, then $x^2 = 0$, hence $x^2 \geq 0$.

If $x \neq 0$, then by Remark 7.25, we have that $x^2>0$. 

\renewcommand\qedsymbol{QED}
\end{proof}


\end{lemma}

\begin{corollary}
	$0 < 1$.
        
\begin{proof}
We know that $0 \neq 1$, and $1=1^2$, hence $1>0$.

\renewcommand\qedsymbol{QED}
\end{proof}


\end{corollary}

\begin{theorem}
	If $F$ is an ordered field, then $F$ has no first or last point.  
        
\begin{proof}
We know that $0<1$, hence $0+(-1)<1+(-1)$, hence $-1<0$. 

Let $x \in F$ be any element of the ordered field. Then, $x+(-1)<x+0<x+1$, hence $x+(-1)<x<x+1$.

\renewcommand\qedsymbol{QED}
\end{proof}


\end{theorem}

\begin{theorem}
	
	The rational numbers $\bbQ$ form an ordered field.
\begin{proof}
First, we want to show that addition respects the ordering.

Let $[\frac{a}{b}],[\frac{c}{d}],[\frac{x}{y}]$ be arbitrary elements of $\bbQ$ with $b,d,y>0$ such that $[\frac{a}{b}]<[\frac{c}{d}]$. 

Since $[\frac{a}{b}]<[\frac{c}{d}]$, we have that $ad<bc$. It follows that
        \begin{align*}
            ad &< bc\\
            adyy &< bcyy\\
            adyy+bdxy &< bcyy+bdxy\\
            aydy+bxdy &< bycy+bydx\\
            (ay+bx)(dy) &< (by)(cy+dx)\\
            [\frac{ay+bx}{by}] &< [\frac{cy+dx}{dy}]\\
            [\frac{a}{b}]+[\frac{x}{y}] &< [\frac{c}{d}] + [\frac{x}{y}].
        \end{align*}

Next, we want to show that multiplication respects the ordering.

Let $[\frac{0}{1}]<[\frac{a}{b}]$ and let $[\frac{0}{1}]< [\frac{c}{d}]$, with $b,d>0$.

Then, we have that $0\cdot b<1\cdot a$ and $0\cdot d<1\cdot c$, hence $0<a$ and $0<c$, hence $0<a\cdot c$. Thus, $0\cdot b\cdot d<1 \cdot a \cdot c$, thus, $[\frac{0}{1}]<[\frac{a\cdot c}{b \cdot d}]$. Therefore, $[\frac{0}{1}]<[\frac{a}{b}] \cdot [\frac{c}{d}]$.


\renewcommand\qedsymbol{QED}
\end{proof}
    
\end{theorem}


There is an important sense in which every ordered field contains a ``copy" of the rational numbers. Let $F$ be an ordered field with additive identity $0_F$ and multiplicative identity $1_F$. To prevent ambiguity, denote the additive identity in the rational numbers by $0_\bbQ$ and the multiplicative identity in the rational numbers by $1_\bbQ$.\medskip



\begin{theorem}
	Any ordered field $F$ contains  a copy of the rational numbers $\bbQ$ in the following sense.  There exists an injective map $i\colon \bbQ \arr F$ that respects all of the axioms for an ordered field. In particular:
	\begin{itemize}
		\item $i(0_\bbQ) = 0_F$.
		\item $i(1_\bbQ) = 1_F$.
		\item If $a, b\in \bbQ$, then $i(a+b) = i(a) + i(b)$.
		\item If $a, b\in \bbQ$, then $i(a\cdot b)=i(a)\cdot i(b)$.
		\item If $a, b\in \bbQ$ and $a < b$, then $i(a)< i(b)$.
	\end{itemize}
	
	    Hint: start by defining a function $j\colon\bbZ\arr F$ satisfying the properties above, then use
    $j$ to define $i$.


\begin{proof}

Define $j$ inductively by $j(0)=0, j(n)=j(n-1)+1$ and $j(-n)=-j(n)$ for all $n \in \bbN$.

Now, we want to check that the above properties, along with the property that $j(-z)=-j(z)$, hold.

\begin{enumerate}
    

\item $j(0)=0$ and we are done.

\item $j(1)=j(0)+1=0+1=1$ and we are done.

\item $j(-n)=-j(n)$ for all $n\in \bbN$.

$0=j(-0)=-j(0)$.

$j(-(-n))=j(n)$, hence $-j(-(-n))=-j(n)=j(-n)$, hence $j(-(-n)) = -j(-n)=j(n)$.

\item For addition, case 1 is that $a\geq 0, b \geq 0$, then we want to show that $j(a+b)=j(a)=j(b)$ by induction in $a$. 

If $a=0$, then $j(0+b)=j(b)=0+j(b)=j(0)+j(b)=j(a)+j(b)$.

If $a \geq 1$, then $j(a+b)=j((a-1)+(b+1))=j(a-1)+j(b+1)=j(a-1)+j(b)+1=j(a-1)+1+j(b)=j(a)+j(b)$.

Case 2 is that $a <0, b \geq 0$. Then, $a=-n$ for some $n \in \bbN$. We want to show that $j(-n+b)=j(-n)+j(b)$ by induction on $n$. 

If $n=1$, then $j(-1+b)=j(b-1)=(-1)+j(b)=j(-1)+j(b)$. 

If $n \geq 2$, then $j(-n+b)=j(b-n)=j(-1+b-n)+1=j(-1)+j(b)+j(-n)+1=j(-1)+j(b)+j(-n)+j(1)=j(b)+j(n)$.

Case 3 is that $b <0, a \geq 0$. This is analogous to case 2.

Case 4 is that $a<0,b<0$. Then $a=-p, b=-q$ for some $p,q \in \bbN$. We want to show that $j(-p-q)=j(-p)+j(-q)$. We have that $j(-p-q)=-j(p+q)=-j(p)-j(q)=j(-p)+j(-q)$.

\item For multiplication, case 1 is that $a\geq 0, b\geq 0$. We want to show that $j(a \cdot b)=j(a)\cdot j(b)$ by induction on $a$. 

For $a=0$, we have that $j(a\cdot b)=j(0)=0=j(0)\cdot j(b)$.

For $a \geq 1$, we have that $j(a \cdot b)=j(((a-1)+1)\cdot b)=j((a-1)\cdot b + b)=j((a-1)\cdot b)+j(b)=j(a-1)\cdot j(b)+j(b)=j((a-1)+1)\cdot j(b)=(j(a-1)+j(1))\cdot j(b) = j(a-1+1)\cdot j(b)=j(a) \cdot j(b)$.

Case 2 is that $a<0, b \geq 0$. Then, $a=-n$ for some $n \in \bbN$. Then, we want to show that $j(-n \cdot b) = j(-n) \cdot j(b)$. We have that $j(-n \cdot b)=-j(n \cdot b)=-(j(n) \cdot j(b))=-j(n) \cdot j(b)=j(-n) \cdot j(b)$. 

Case 3 is that $a\geq 0, b < 0$. This is analogous to case 2.

Case 4 is that $a<0,b<0$. Then, $j(a\cdot b)=j(-(-(a\cdot b)))=-j((-a)\cdot b)=-j((-a)\cdot j(b))=(-j(-a))\cdot j(b)=j(a) \cdot j(b)$.

\item We want to show order-preserving property.

If $a<b$, then there exists $n \in \bbN$ such that $a+n=b$. Then $j(a)+j(n)=j(b)$. Thus, it suffices to show that $j(n)>0$. We will show this by induction on $n$.

If $n=1$, then $j(1)=1>0$. 

If $n\geq 2$, then $j(n)=j((n-1)+1)=j(n-1)+j(1)>0$, where the last step follows since $j(n-1)>0$ (by inductive hypothesis) and $j(1)>0$ (by definition of $j$).

\end{enumerate}

Now, we proceed to define $i \colon \bbQ \rightarrow F$ as $i(\frac{a}{b})=\frac{j(a)}{j(b)}$.

First, we want to show that $i$ is well-defined. Since $b \neq 0$, we have that $j(b) \neq 0$. Suppose $\frac{a}{b}=\frac{c}{d}$, then we have that $ad=cb$, hence $j(ad)=j(cb)$, hence $j(a)\cdot j(b)=j(c) \cdot j(b)$, hence $\frac{j(a)}{j(b)}=\frac{j(c)}{j(d)}$.

Next, we want to check that all the desired properties hold.

\begin{enumerate}
\item 
$i(0)=\frac{j(0)}{j(1)}=j(0)=0$.

\item 
$i(1)=\frac{j(1)}{j(1)}=j(1)=1$.

\item 
$i(\frac{a}{b}+\frac{c}{d})=i(\frac{ad+bc}{bd})= \frac{j(ad+bc)}{j(bd)}=\frac{j(ad)+j(bc)}{j(bd)}=\frac{j(a)}{j(b)}+\frac{j(c)}{j(d)}=i(\frac{a}{b}) + i(\frac{c}{d})$. 

\item 
$i(\frac{a}{b}\cdot \frac{c}{d})=i(\frac{ac}{bd})= \frac{j(ac)}{j(bd)}=\frac{j(a)\cdot j(c)}{j(b)\cdot j(d)}=\frac{j(a)}{j(b)}\cdot \frac{j(c)}{j(d)}=i(\frac{a}{b}) \cdot i(\frac{c}{d})$. 

\item 
If $\frac{a}{b}<\frac{c}{d}$, then $ad<cb$, hence $j(ad)<j(cb)$, hence $j(a) \cdot j(d) < j(c) \cdot j(b)$, hence $\frac{j(a)}{j(b)}<\frac{j(c)}{j(d)}$, hence $i(\frac{a}{b})<i(\frac{c}{d})$.

This completes the proof.

\end{enumerate}




\renewcommand\qedsymbol{QED}
\end{proof}

\end{theorem}


Since we know that $\bbQ$ is an ordered field, this result says that $\bbQ$ is the ``minimal'' ordered field.\bigskip


\begin{center}
{\em Appendix: The Real Numbers are an Ordered Field}
\end{center} 
This appendix is concerned with proving that the continuum $\bbR$ is an ordered field.  Addition and multiplication on $\bbR$ are defined in terms of addition and multiplication on $\bbQ$, so we will use $\oplus$ and $\otimes$ for addition and multiplication of real numbers, to make sure that there is no confusion with $+$ and $\cdot$ on $\bbQ$.



\begin{definition}
	We define $\oplus $ on $\bbR$ as follows. Let $A,B\in\bbR$ be Dedekind cuts. Define
	\begin{align*}
		A \oplus B = \{a + b \mid a \in A\text{ and }b \in B\}.
	\end{align*}
\end{definition}

\begin{exercise}
	\begin{enumerate}[(a)]
		\item Prove that $A \oplus B$ is a Dedekind cut.        
\begin{proof}
First, $A \neq \emptyset$, hence there exists $a \in A$ and likewise $B \neq \emptyset$, hence there exists $b \in B$, hence there exists $a+b \in A \oplus B$. Hence $A \oplus B \neq \emptyset$. 

Also, by trichotomy, either $A<B$ or $A>B$ or $A=B$, hence without loss of generality suppose $A \subseteq B$. There exists $c \in \bbQ \setminus B$, hence $2c \notin A \oplus B$, hence $A \oplus B \neq \bbQ$.

Second, if $a+b \in A \oplus B$, then we want to show that $a+b-m \in A \oplus B$, for $m \in \bbQ, m>0$. We can write $a+b-m=(a-m)+b$, then $a-m \in A$ since $A$ is closed downwards (it is a Dedekind cut), hence $a+b-m \in A \oplus B$.

Suppose $a+b \in A \oplus B$. Then, since $A$ has no last point, there exists $a' \in A$ such that $a'>a$. Then, $a'+b \in A \oplus B$ and $a'+b>a+b$. Hence $A \oplus B$ has no last point.

Therefore $A \oplus B$ is a Dedekind cut.


\renewcommand\qedsymbol{QED}
\end{proof}

		\item Prove that $\oplus$ is commutative and associative. 

\begin{proof}
$A \oplus B = \{a + b \mid a \in A\text{ and }b \in B\} = \{b + a \mid a \in A\text{ and }b \in B\} = B \oplus A$, since $A \subset \bbQ$ and $B \subset \bbQ$ and addition on $\bbQ$ is commutative.

$(A \oplus B)\oplus C = \{(a + b)+c \mid a \in A\text{ and }b \in B \text{ and }c \in C\} = \{a + (b+c) \mid a \in A\text{ and }b \in B \text{ and }c \in C\} = A \oplus (B \oplus C)$, since $A \subset \bbQ$ and $B \subset \bbQ$ and $C \subset \bbQ$, and addition on $\bbQ$ is associative.

\renewcommand\qedsymbol{QED}
\end{proof}

		\item Prove that if $A \in \bbR$ then $A = \mathbf{0} \oplus A$. 
                
\begin{proof}
$\mathbf{0}=i(0)=D_0=\{q \in \bbQ \mid q<0\}$.

Let $m \in D_0$. Then, $m<0$, hence for all $a \in A, a+m<a$. Since $A$ is closed downwards, $a\in A$ implies $a+m \in A$, hence $\mathbf{0} \oplus A=A$.

\renewcommand\qedsymbol{QED}
\end{proof}

\end{enumerate}
\end{exercise}

\begin{lemma}
	\label{lem:in A close to not in A}
	If $A\in\bbR$, $r\in\bbQ$, and $r>0$, then there exists $s\in A$ such that $s+r\not\in A$.
        
\begin{proof}
Suppose there does not exist $s \in A$ such that $s+r \notin A$. Then, for all $s\in A, s+r \in A$ and by the property of being closed downwards, if $q \in \bbQ$ such that $q<s+r$ then $q \in A$. Since $s+r \in A$ for all $s \in A$, we have that $A$ has no upper bound. This implies that $A =\bbQ$ and this is a contradiction.
 
\renewcommand\qedsymbol{QED}
\end{proof}
\end{lemma}

\begin{lemma}
	Let $A\in\bbR$. Let
	\begin{align*}
	B= \{r \in \bbQ \mid -r \notin A\text{ but } {-r} \text{ is not a first point of }\bbQ \setminus A\}.
	\end{align*}
	Then,
	\begin{enumerate}[(a)]
		\item $B \in \bbR$.
                
\begin{proof}
If $A = \mathbf{0}$, then $A=B$ and we are done. (This is because, for all $a \in A, -a \notin A$, and $0$ is the first point of $\bbQ \setminus A$, hence $-a$ is not a first point of $\bbQ \setminus A$).

Case 1: $A < \mathbf{0}$. Then there exists $q \in \mathbf{0}$ such that $q \notin A$. Then $q<-q$, hence $-q \notin A$ and $-q$ is not a first point of $\bbQ \setminus A$, hence $q \in B$, hence $B \neq \emptyset$.

Let $a \in A$, then $a<0$, hence $-a \in \bbQ \setminus B$, since $-(-a)=a \in A$. Hence, $B \neq \bbQ$.

Let $b \in B$, and let $b' \in \bbQ$ be such that $b'<b$. Then, we have that $-b'>-b$, hence $-b' \notin A$ (since $-b \notin A$ which implies $-b$ is an upper bound of $A$, and $-b'>-b$  which implies $-b'$ is an upper bound of $A$). Also, $-b$ is not a first point of $\bbQ \setminus A$, and $-b'>-b$, hence $-b'$ is not a first point of $\bbQ \setminus A$. Thus, $b' \in B$.

Let $b \in B$, then we have that $-b \notin A$ and $-b$ is not a first point of $\bbQ \setminus A$. Since $\bbQ$ is dense, there exists $c \in B$ such that $c<-b$, $c \notin A$ and $c$ is not a first point of $\bbQ \setminus A$. Then, $-c\in B, -c>b$, hence $B$ has no last point.

Case 2: $A>\mathbf{0}$. We know that $\bbQ \setminus A$ is non-empty, hence let $d \in \bbQ \setminus A$ such that $d$ is not a first point of $\bbQ \setminus A$. Then, $-d \in B$, hence $B \neq \emptyset$.

Let $q \in A \setminus \mathbf{0}$, then $-q<q$ and since $A$ is closed downwards, $-q \in A$, hence $q \notin B$, hence $B \neq \bbQ$.

Let $b \in B$, and let $b' \in \bbQ$ be such that $b'<b$. Then, $-b'>-b$, and since $-b \notin A$ ($-b$ is an upper bound of $A$) and $-b$ is not a first point of $\bbQ \setminus A$, by $-b'>-b$ we have that $-b' \notin A$ and $-b'$ is not a first point of $\bbQ \setminus A$. Hence, $b' \in B$. 

Let $b \in B$, then $-b \notin A$ and $-b$ is not a first point of $\bbQ \setminus A$. Since $\bbQ$ is dense, there exists $c \in \bbQ \setminus A$ such that $c<-b$, $c \notin A$ and $c$ is not a first point of $\bbQ \setminus A$. Then, $-c\in B, -c>b$, hence $B$ has no last point.

This demonstrates that $B$ is a Dedekind cut, i.e., $B \in \bbR$.

\renewcommand\qedsymbol{QED}
\end{proof}
		\item $A \oplus B = \mathbf{0}$.
            
\begin{proof}
For the first containment, let $m \in A \oplus B$, then $m=a+b$ for some $a\in A, b\in B$. We have that $-b \notin A$, hence $-b$ is an upper bound of $A$, hence $a<-b$. Then, $a+b<-b+b=0$, hence $m \in \mathbf{0}$. Thus, $A \oplus B \subseteq \mathbf{0}$.

For the second containment, we want to show that if $m<0$, then there exist $a\in A, b\in B$ such that $m=a+b$. 

Since $m<0, -m>0$. Hence, there exists $a \in A$ such that $a+(-m) \notin A$. Since $A$ has no last point, there exists $a' \in A$ such that $a'>a$ and $a'+(-m) \notin A$. 

Then, $a'+(-m)>a+(-m)$. From this we know that $a'+(-m)$ is not the first point of $\bbQ \setminus A$. Hence, $-(a'+(-m)) \in B$.

Then, $a'+(-(a'+(-m))) = a' - (a' -m) = m$, hence $m \in A \oplus B$. Thus, $\mathbf{0} \subseteq A \oplus B$.

Therefore, $A \oplus B = \mathbf{0}$. This completes the proof.

\renewcommand\qedsymbol{QED}
\end{proof}
		
	\end{enumerate}

\end{lemma}
\begin{theorem}
	Every $A \in \bbR$ has an additive inverse $-A$ such that $A \oplus (-A) = \mathbf{0}$. 
        
\begin{proof}
Consider the construction of $B$ above and set $B =-A$, then we are done.

\renewcommand\qedsymbol{QED}
\end{proof}
\end{theorem}

\begin{exercise}
	Show that $A < \mathbf{0} \Longleftrightarrow \mathbf{0} < -A$.
        
\begin{proof}
First, we want to show that $\mathbf{0} \subset -A$. Let $x \in \mathbf{0}$, then $-x>0$ hence $-x \notin A$. Also, $0<-x$ and $0\in \bbQ \setminus A$, hence $-x$ is not the first point of $\bbQ \setminus A$. Hence, $x \in -A$. 


Second, we want to show that $-A \setminus \mathbf{0} \neq \emptyset$. We know that $A \subsetneq \mathbf{0}$, hence there exists $b \in \mathbf{0} \setminus A$. Since $\bbQ$ is dense, we can find $b' \in \mathbf{0} \setminus A$ such that $b' \notin A$ and $b'>b$ which implies that $b'$ is not the first point of $\bbQ \setminus A$. Then, $-b' \in -A$ and since $b'<0$, we have that $-b' > 0$. Hence, $-b' \in -A \setminus \mathbf{0}$. 

This completes the proof.

\renewcommand\qedsymbol{QED}
\end{proof}
\end{exercise}

\begin{exercise}
	Show that if $p\in\bbQ$ and $A=\{x\in\bbQ\mid x<p\}$, then $-A=\{x\in\bbQ\mid x<-p\}$.
        
\begin{proof}
For the first containment, if $a' \in -A$, then $-a'\notin A$, hence $-a'>p$, hence $a'<-p$.

For the second containment, if $a<-p$, then $-a>p$, hence $-a \notin A$, hence $a \in -A$.

(We note that the first point of $\bbQ \setminus A$ is $p$).

\renewcommand\qedsymbol{QED}
\end{proof}
\end{exercise}



\begin{exercise}
	Show that $<$ satisfies additivity; i.e., if  $A,B,C\in\bbR$ with $A<B$ then $A\oplus C<B\oplus C$. 
        
\begin{proof}
We have that $A<B$, hence there exists $b \in B$ such that $b \notin A$, then $b\geq a$ for all $a \in A$. Since $B$ has no last point, there exists $b' \in B$ such that $b'>b$, hence $b'>a$ for all $a \in A$.

Then, $b'+c\in B \oplus C$, but $b'+c>a+c$ for all $a \in A, c \in C$, hence $b'+c \notin A \oplus C$.

Hence, $A \oplus C < B \oplus C$.

\renewcommand\qedsymbol{QED}
\end{proof}
\end{exercise}


\begin{definition}
	For $A,B\in\bbR$, $\mathbf{0}<A$, $\mathbf{0}<B$, we define
	\begin{align*}
		A \otimes B = \{r \in \bbQ \mid r \leq 0\} \cup \{ab \mid a\in A, b\in B, a > 0, b > 0\}.
	\end{align*}
	If $A = \mathbf{0}$ or $B = \mathbf{0}$ we define $A \otimes B = \mathbf{0}$.
	If $A<\mathbf{0}$ but $\mathbf{0}<B$ we replace $A$ with $-A$ and use the definition of multiplication of positive elements. Hence, in this case,
	\begin{align*}
		A\otimes B=-[(-A)\otimes B].
	\end{align*}
	Similarly, if $\mathbf{0}<A$ but $B<\mathbf{0}$, then 
	\begin{align*}
		A\otimes B=-[A\otimes (-B)],
	\end{align*}
	and 
	if $A<\mathbf{0}, B<\mathbf{0}$ then
	\begin{align*}
		A \otimes B= (-A) \otimes (-B).
	\end{align*}
\end{definition}

\begin{exercise}
	\begin{enumerate}[(a)]
		\item Show that if $A, B \in \bbR$, then $A \otimes B \in \bbR$.
        
\begin{proof}
$-1 \in \{r \in \bbQ \mid r \leq 0\}$, hence $-1 \in A \otimes B$, hence $A \otimes B \neq \emptyset$.

$A\neq \bbQ$ and $B \neq \bbQ$, hence there exist $a',b' \in \bbQ$ such that $a'$ is an upper bound of $A$ and $b'$ is an upper bound of $B$. Then, $a'b'$ is an upper bound of $A \otimes B$, hence $a'b' \notin A \otimes B$, hence $A \otimes B \neq \bbQ$.

Suppose $s \in A \otimes B$ and $t\in \bbQ$ such that $t<s$, then we want to show that $t \in A\otimes B$.

Case 1: if $t\leq0$, then $t\in A \otimes B$.

Case 2: if $t>0$, then $s>0$, hence $s=ab$ for $a\in A, b\in B$. Then, $t=s-q=ab-q=b(a-\frac{q}{b})$ such that $q\in \bbQ, q >0$. Since $A$ is closed downwards, we have that $(a-\frac{q}{b})\in A$, hence $b(a-\frac{q}{b})=ab-q=s-q=t \in A\otimes B$. This shows that $A\otimes B$ is closed downwards.

Suppose $m \in A \otimes B$, then we want to show that there exists $m' \in A\otimes B$ such that $m'>m$.

Case 1: $m\leq 0$. Since $A,B$ have no last point, there exists $x \in A \otimes B$ such that $x>m$. 

Case 2: $m>0$. Then, $m=ab$ for some $a\in A, b\in B$. Since $A$ has no last point, there exists $a'\in A$ such that $a'>a$, then $m'=a'b>ab=m$. 

\renewcommand\qedsymbol{QED}
\end{proof}

		\item Show that $\otimes$ is commutative and associative. 
		        
\begin{proof}
First, we want to show that $A\otimes B=B \otimes A$. 

We have that \begin{align*}
            A\otimes B &= \{r\in\bbQ\mid r\leq 0\}\cup\{ab\mid a\in A,b\in B,a>0,b>0\}\\
            &= \{r\in\bbQ\mid r\leq 0\}\cup\{ba\mid b\in B,a\in A,b>0,a>0\}\\
            &= B\otimes A.
        \end{align*}

Next, we want to show that $(A\otimes B)\otimes C=A \otimes (B \otimes C)$. Let $x \in (A\otimes B)\otimes C$. 

Case 1: $x \leq 0$. Then, $x \in \{r\in\bbQ\mid r\leq 0\}$, hence $x \in A \otimes (B \otimes C)$. 

Case 2: $x>0$. Then, $x=dc$, where $d\in A \otimes B, c\in C$, $d>0,c>0$. Further, $d=ab$ for $a\in A,b\in B$ such that $a,b>0$. Let $y=bc$, then $y \in B\otimes C$. Then, $x=ay$, hence $x\in A \otimes (B\otimes C)$. 

This completes the first containment. The second containment is analogous.

\renewcommand\qedsymbol{QED}
\end{proof}

        \item Show that if $A, B \in \bbR$, $\mathbf{0} < A$, and $\mathbf{0} < B$, then $\mathbf{0} < A\otimes B$. 
		        
\begin{proof}
We have that $0 \subsetneq A, 0 \subsetneq B$, hence there exist $a\in A,b\in B$ such that $a,b>0$. Then $ab>0$ and $ab\in A \otimes B$, hence $\mathbf{0}<A \otimes B$.

\renewcommand\qedsymbol{QED}
\end{proof}

        \item Let $\mathbf{1} = \{x \in \bbQ \mid x < 1\}$. Show that if $A \in \bbR$, then $\mathbf{1} \otimes A = A$. 
	        
\begin{proof}
Let $a \in A$, then let $a'\in A$ such that $a'=a+\epsilon$ and $\epsilon << a'$. Then, \[
a=a'-\epsilon = a'(1-\frac{\epsilon}{a'}) = a'(\frac{a'-\epsilon}{a'}).
\]

Then, we have that $a' \in A$ by definition, and $\frac{a'-\epsilon}{a'} \in \mathbf{1}$ such that $\frac{a'-\epsilon}{a'}>0$. This completes the first containment, namely $A \subset \mathbf{1} \otimes A$.

Let $a \in \mathbf{1} \otimes A$. Then, $a = jk$, for $j\in \mathbf{1},j>0$ and $k \in A, k>0$. Then, $a<k$ and $k\in A$, which implies $a \in A$ by the property of being closed downwards. This completes the first containment, namely $\mathbf{1} \otimes A \subset A$.

\renewcommand\qedsymbol{QED}
\end{proof}

    \end{enumerate}
\end{exercise}

\begin{lemma}
	If $A \in \bbR$, $A > 0$, $r \in \bbQ$, and $r > 1$, then there exist $s\in A$ such that $rs\not\in A$.
	

	(Hint: Prove that $\{r^n \mid n \in \bbN\}$ is unbounded in $\bbQ$.)

        
\begin{proof}
First, we want to show that $\{r^n \mid n \in \bbN\}$ has no upper bound. Let $r^k \in \{r^n \mid n \in \bbN\}$, then $r^{k+1} \in \{r^n \mid n \in \bbN\}$ and $r^{k+1}>r^k$.

Suppose there does not exist $s \in A$ such that $sr \notin A$. Then, for all $s\in A, sr \in A$ and by the property of being closed downwards, if $q \in \bbQ$ such that $q<sr$ then $q \in A$. Since $sr \in A$ for all $s \in A$, we have that $A$ has no upper bound. This implies that $A =\bbQ$ and this is a contradiction.

\renewcommand\qedsymbol{QED}
\end{proof}

\end{lemma}



\begin{lemma}
	Let $A\in\bbR$ with $\mathbf{0}<A$. Let
	\begin{align*}
		B = \{r \in \bbQ \mid r \leq 0\} \cup \left\{r \in \bbQ \mathrel{}\middle|\mathrel{} r > 0\text{ and }\frac{1}{r} \notin A, \text{ but } \frac{1}{r} \text{ is not a first point of }\bbQ \setminus A\right\}.
	\end{align*}
	Then,
	\begin{enumerate}[(a)]
		\item $B \in \bbR$.
                
\begin{proof}
$-1\in \{r \in \bbQ \mid r \leq 0\}$, hence $-1\in B$, hence $B \neq \emptyset$.

There exists $a\in A$ such that $a>0$. Then, there exists $n\in \bbN$ such that $0<\frac{1}{n}<a$. Hence, $\frac{1}{n} \in A$ by the property of being closed downwards. Then, $[\frac{n}{1}] \notin B$, hence $B \neq \bbQ$.

Suppose $b \in B$ and $c<b$. If $c \leq 0$, then $c \in B$ and we are done. If $c>0$, then we have that $\frac{1}{c}>\frac{1}{b}$, $\frac{1}{b} \notin A$ and $\frac{1}{b}$ is not a first point of $\bbQ \setminus A$. From this it follows that $\frac{1}{c} \notin A$ and $\frac{1}{c}$ is not a first point of $\bbQ \setminus A$. Hence, $c \in B$.

Let $b \in B$. If $b\leq 0$ then we know there exists $b'\in B$ such that $b'>b$ and we are done. 

If $b>0$, then $\frac{1}{b} \notin A$ and $\frac{1}{b}$ is not a first point of $\bbQ \setminus A$. Suppose $p$ is the first point of $\bbQ \setminus A$. Then, by the Archimedean property, there exists $c \in \bbN$ such that $p<\frac{1}{c}<\frac{1}{b}$, hence $c>b$ and $c\in B$. 

This proves that $B \in \bbR$.



\renewcommand\qedsymbol{QED}
\end{proof}

		\item $A \otimes B = \mathbf{1}$.
	        
\begin{proof}
Let $a\in A, b\in B$. Then, we know that $\frac{1}{b} \notin A$ and $\frac{1}{b}$ is not a first point of $\bbQ \setminus A$. Hence, for all $a\in A$, we have that $a< \frac{1}{b}$, hence $ab<1$. This completes the first containment, namely $A \otimes B \subset \mathbf{1}$.

Let $m \in \mathbf{1}$. If $m\leq 0$ then $m \in A\otimes B$ and we are done. If $m>0$, that is if $0<m<1$, then since $\mathbf{0}<A$, there exists $a\in A$ such that $a>0$. Then, $ab>0$ and by the above reasoning, $ab<1$. Hence, we can set $m=ab$, that is, we know that $m$ can be written as $ab$ for some $a\in A, b\in B$. This completes the first containment, namely $\mathbf{1} \subset A \otimes B$.

This completes the proof.


\renewcommand\qedsymbol{QED}
\end{proof}

    \end{enumerate}
\end{lemma}

\begin{lemma}
	Every $A \in \bbR$ with $\mathbf{0} < A$ has a multiplicative inverse $A^{-1}$ such that $A \otimes A^{-1} = \mathbf{1}$.
        
\begin{proof}
Consider the above construction of $B$ and set $B=A^{-1}$, then we are done.

\renewcommand\qedsymbol{QED}
\end{proof}

\end{lemma}

\begin{lemma}
	Let $A \in \bbR$ with $A < \mathbf{0}$. Then, $(-A)^{-1}$ has already been defined, and $A \otimes (-(-A)^{-1}) = \mathbf{1}$.
        
\begin{proof}
If $A<0$ then $-A>0$, hence $(-A)^{-1}$ has already been defined above. We have that $(-A)\otimes (-A)^{-1}=\mathbf{1}=A\otimes (-(-A)^{-1})$ by commutativity of $\otimes$.


\renewcommand\qedsymbol{QED}
\end{proof}

\end{lemma}

\begin{theorem}
	Every $A \in \bbR$ with $A \neq \mathbf{0}$ has a multiplicative inverse $A^{-1}$ such that $A \otimes A^{-1} = \mathbf{1}$. 
        
\begin{proof}
Since $A\neq \mathbf{0}$, either $A<\mathbf{0}$ or $A>\mathbf{0}$. 

If $A<\mathbf{0}$ then consider Lemma 7.44 and we are done.

If $A > \mathbf{0}$ then consider Lemma 7.43 and we are done.

\renewcommand\qedsymbol{QED}
\end{proof}

\end{theorem}

\begin{exercise}
	Show that if $p \in \bbQ$ with $p \neq 0$ and $A= \{x \in \bbQ \mid x < p\}$, then $A^{-1} = \{x \in \bbQ \mid x < \frac{1}{p}\}$.
        
\begin{proof}
Let $m \in \{x \in \bbQ \mid x < \frac{1}{p}\}$, then $m<\frac{1}{p}$. Hence, $\frac{1}{m}>p$, hence $\frac{1}{m} \notin A$. We also know that $p$ is the first point of $\bbQ \setminus A$, and since $\frac{1}{m}>p$, we have that $\frac{1}{m}$ is not the first point of $\bbQ \setminus A$. Hence, $m \in A^{-1}$. This completes the first containment, namely $\{x \in \bbQ \mid x < \frac{1}{p}\} \subset A^{-1}$.

Let $m \in A^{-1}$. Then, $\frac{1}{m} \notin A$ and $\frac{1}{m}$ is not the first point of $\bbQ \setminus A$, hence $\frac{1}{m}>p$. Hence $m<\frac{1}{p}$, hence $m \in \{x \in \bbQ \mid x < \frac{1}{p}\}$. This completes the second containment, namely $A^{-1} \subset \{x \in \bbQ \mid x < \frac{1}{p}\}$.

This completes the proof.

\renewcommand\qedsymbol{QED}
\end{proof}

\end{exercise}

\begin{theorem}
	$\bbR$ is an ordered field. 
        
\begin{proof}
First, we want to show that $\oplus$ respects the ordering. Suppose $A<B$, then $A\oplus C = \{r\in \bbQ \mid r\leq 0\} \cup \{a+c \mid a\in A, c\in C, a>0, c>0\}$. And similarly, $B\oplus C = \{r\in \bbQ \mid r\leq 0\} \cup \{b+c \mid b\in B, c\in C, b>0,c>0\}$. 

Since $A<B$, we know that for all $a \in A, a\in B$. Hence, for all $d \in A\oplus C, d \in B \oplus C$. However, there exists $b \in B$ such that $b \notin A$, i.e., $b$ is an upper bound of $A$. Then, $a\leq b$ for all $a \in A$. Then, there exists $b' \in B$ such that $b'>b$, hence $a< b'$ for all $a\in A$. Hence $a+c<b'+c$, hence $A\oplus C < B \oplus C$.

Second, we want to show that $\otimes$ respects the ordering.

Suppose $A,B \in \bbR$ such that $A>\mathbf{0},B>\mathbf{0}$. Then, there exist $a\in A$ and $b\in B$ such that $a>0,b>0$, hence $ab>0$. Hence, $A \otimes B >\ \mathbf{0}$.

This completes the proof.

\renewcommand\qedsymbol{QED}
\end{proof}

\end{theorem}
\newpage

\begin{center}
{\em Additional Exercises}
\end{center}

{\em In these exercises you should assume the results from the Appendix, whether or not you proved these in class. Also, the function $i:\bbQ\to\bbR$ referred to in exercises 1 and 2 is the function defined in Script 6 - so $i(q)=\{x\in\bbQ\mid x<q\}.$}

\begin{enumerate}

\item
\begin{enumerate}
\item[a)] Let $q\in\bbQ.$
Prove that  $i(q) \oplus i(-q)=\bf{0.}$ 
\item[b)] Show that for $q_1,q_2\in \bbQ, $
 $$i(q_1)\otimes i(q_2)=i(q_1 q_2).$$

 \end{enumerate}







\item 
\begin{enumerate}
\item[a)] Define $\sqrt{2}=\{x\in\bbQ\mid x< 0\}\cup\{x\in \bbQ\mid x^2<2\}.$ We saw in class that $\sqrt{2}\in\bbR.$\\Show that
$\sqrt{2}$ is not a rational cut, i.e. there is no $p\in\bbQ$ such that $\sqrt{2}=\{x\in\bbQ\mid x<p\}.$
\item[b)] Show that $\sqrt{2} \otimes i(x)\in \bbR\setminus i(\bbQ),$ for $x\in\bbQ, x>0.$ (Note that in fact this holds for all $x\neq 0,$ but you need only show it for $x>0.$)
\item[c)] Show that the irrationals are dense in the reals - more precisely, show that $\bbR\setminus i(\bbQ) $ is dense in $\bbR.$
\item[d)] {\em (Challenge)} Show that $\sqrt{2}\otimes \sqrt{2}=\{x\in\bbQ\mid x<2\}.$ 
\end{enumerate}

\item Let $A \subset \bbR$ and assume that $\sup A$ exists.  Define $-A := \{ x \in \bbR : -x \in A\}$.  Is it true that $-\sup A = \sup -A$?

\item 
\begin{enumerate}
		\item Let $A,B$ be bounded subsets of $\bbR.$
		
			Define $A+B = \{c \in \bbR: c = a+b$ for some $a\in A, b\in B\}$.  Is it true that $\sup A + \sup B = \sup(A+B)$?
		\item Let $f,g: \bbR \to \bbR$ be functions.  Is it true that $\sup f(\bbR) + \sup g(\bbR)= \sup (f+g)(\bbR)$?  (Note: here we define $f+g:\bbR \to \bbR$ to be the function $(f+g)(c) = f(c) + g(c)$ for all $c\in \bbR$.)
	\end{enumerate}



\end{enumerate}
\end{document}

